\section{Negative results and limitations}
\label{sec:limits}

We keep negative results, failed runs and reruns in the record. This section collects them.

\subsection{Negative results}

\paragraph{Training increased susceptibility to wrong-answer injection.}
The pooled attack success of a sentence naming a wrong answer rose from 21.0\% to 39.4\%
(Section~\ref{sec:injection}), even though training included an injection-detection task. This
is the most important open problem for the tasks Bobcat targets, several of which read
untrusted input. A structurally closed output does not help here.

\paragraph{The residual pointer head was not adopted.}
In joint training the pointer's scale stayed at zero (Section~\ref{sec:arms}). We therefore
retested it with the backbone and any adapter frozen and run without gradients, training only
the head for one epoch (1,592 steps of 32 questions, head learning rate $10^{-3}$) with the
scale initialized at 1.0 (Table~\ref{tab:pointer}). On the frozen base model the pointer adds
about one point (+1.1 [+0.1, +2.1] over zero-shot), against +7.9 for LoRA, and it reduced its own
scale from 1.0 to 0.016. On top of \expone{} it costs 2.1 points ([$-2.9$, $-1.2$]) and makes
the model strongly overconfident (NLL 0.770; fitted temperature 3.96). In both cases it hurts
255-candidate title selection (base 65.3\% to 60.0\%; \expone{} 81.3\% to 57.3\%). Our
explanation is that the training data contain few long candidate lists, so a bilinear score
between candidate end states does not generalize to long lists. Retesting a candidate
comparison structure first needs training data with long lists and a normalization that is
invariant to the number of candidates.

\begin{table}[t]
\centering
\small
\caption{Pointer head retest on the development split (head-only training, scale initialized
at 1.0). $\Delta$ in points with 95\% paired component-bootstrap intervals; the reference arm
is named in each row.}
\label{tab:pointer}
\setlength{\tabcolsep}{3.5pt}
\begin{adjustbox}{max width=\linewidth}
\begin{tabular}{@{}lrllrrrrrrr@{}}
\toprule
Arm & Macro & $\Delta$ [95\%] & Against & Tool & NLL & $T$ & ECE@$T$ & Monitor & K255 & Scale \\
\midrule
Zero-shot (reference) & 85.5 & -- & -- & 81.8 & 0.526 & 1.26 & 0.045 & 83.1 & 65.3 & -- \\
Pointer only, base frozen & 86.6 & +1.1 [+0.1, +2.1] & zero-shot & 80.7 & 0.486 & 1.40 & 0.028 & 87.6 & 60.0 & 0.016 \\
\expone{} (LoRA SFT) & 93.4 & +7.9 [+6.8, +9.0] & zero-shot & 95.0 & 0.236 & 1.15 & 0.011 & 95.5 & 81.3 & -- \\
\expone{} + pointer, frozen & 91.3 & $-2.1$ [$-2.9$, $-1.2$] & \expone{} & 93.4 & 0.770 & 3.96 & 0.037 & 94.5 & 57.3 & 0.224 \\
\bottomrule
\end{tabular}
\end{adjustbox}
\end{table}

\paragraph{Reinforcement learning gave no accuracy gain.}
On the student, RL, Brier and CE continuations are indistinguishable in accuracy
(Table~\ref{tab:continuation}); on the GLM reference, proper-score RL and direct Brier tied at
802 of 896 (Section~\ref{sec:glmrl}). This matches the proposition of
Section~\ref{sec:objectives}: with gold labels, the RL estimator carries no information that
the Brier loss lacks. Both results are single-seed. We did not test RL in the setting where
it could add something, with outcome rewards and no gold labels.

\paragraph{Smaller regressions.}
On the ten unseen tool situations, scope judgments were 2.8 points below the untrained model
(71 rows). Title selection with many candidates improved but remains the weakest part of
search (73.7\% at $K=255$ on the final split).

\subsection{Evaluation limitations}

\paragraph{No human review.}
No label of the six-task evaluation has been reviewed by a person, including the 300 rows
that the builder itself flagged as needing review and the classification labels, which we
know to be noisy. The author-written probes and the ten unseen tool situations were written
and labelled by an AI assistant (Claude) and have not been reviewed either; they are sanity
checks.

\paragraph{Language coverage.}
The product evaluation is Korean only. English enters only through 63 probe questions, so we
cannot report per-language results for the product tasks, although our own rules require it.

\paragraph{Few independent units for the held-out task.}
The final tool-call result rests on three situations (96 decisions) and the development
result on six. The component bootstrap cannot capture situation-level variation with so few
clusters.

\paragraph{In-distribution gains.}
Five of the six tasks share generators with the training data (with disjoint components), and
citation verification at 99.8\% (development) and 100.0\% (final) may partly reflect the model
learning the generator's rules. Only the tool-call task measures transfer to an unseen task.

\paragraph{Identifier confound.}
Title questions with 77 or more candidates use different identifiers from those with at most
62 (Section~\ref{sec:identifiers}), so candidate count and identifier choice are confounded in
all $K$ results.

\paragraph{No Score evaluation.}
The product evaluation has no Score primitive. Score behaviour is measured only on the
training-distribution monitor split (normalized MAE) and in ten probe questions.

\subsection{Missing reference and release gates}
\label{sec:gates}

The project's specification defines release gates G0 to G9. The quality reference for G1,
G2 and G9 is the GLM-5.3-Flash direct-logit readout on the same evaluation (baseline B0),
which has not been run. Table~\ref{tab:gates} gives the status of each gate as far as this
paper can speak to it. No gate is claimed as passed.

\begin{table}[t]
\centering
\small
\caption{Status of the specification's release gates. ``Not judged'' means the required
measurement does not exist yet.}
\label{tab:gates}
\setlength{\tabcolsep}{4pt}
\begin{tabular}{@{}lp{4.6cm}p{7.2cm}@{}}
\toprule
Gate & Requirement (abridged) & Status \\
\midrule
G0 & Non-generative path, closed serializer, fixed attack corpus on the release artifact, fault injection & Structural part observed on \expone{} (merged) and the base model with zero violations (Section~\ref{sec:structural}); fault injection tested earlier on the GLM path; the hosted endpoint and network boundary not tested. Semantic attack success reported separately. \\
G1 & Paired 95\% lower bound of the macro difference to GLM direct-logit $> -1$ point & Not judged: reference not run. \\
G2 & G1 in each evaluation language & Not judged: reference not run; no English product evaluation. \\
G3 & ECE $\le 0.05$ per language, NLL and Brier against the reference & Korean final ECE 0.011 at $T$; English not measured; reference not run. \\
G4 & Abstention policy: coverage $\ge 50\%$, selected-error upper bound $\le 2\%$ & Not evaluated. \\
G5 & 77/200/255 candidates without truncation; flip rate $\le 2\%$; permutation TV p95 $\le 0.05$ & 255 candidates handled without truncation; order stability measured with a different statistic (Section~\ref{sec:permutation}); not judged. \\
G6 & Isolation, cache on/off, batch/standalone, resume, with tolerances fixed in advance & Batch and shared-prefix agreement measured (Table~\ref{tab:determinism}); tolerances not fixed in advance; sibling-change and resume tests not run for \expone{}. \\
G7 & W1 warm p50 $\le 150$\,ms, p95 $\le 300$\,ms on a deployment that passed quality & Latency target met on one B300 (97 / 98\,ms HTTP, merged LoRA) and on one H100 with vLLM FP8 (50 / 50\,ms at the engine call); the deployment has not passed the quality gates. \\
G8 & Pre-registered task criteria in both languages, fresh-install reproduction, consistent artifacts & Not met. \\
\bottomrule
\end{tabular}
\end{table}

\subsection{Model size and serving artifact}

The deployed backbone has 27.8B parameters and peaks at 56.5\,GB in serving, above the 1 to
8B dense target of our specification; Gemma-4-12B-it and Qwen3.5-9B remain candidates for the
same training. The merged BF16 model used for the B300 serving measurements differs from the
evaluated unmerged model on 1.4\% of answers (Section~\ref{sec:determinism}), and the merged FP8
artifact of our internal deployment differs on 1.7\% of development answers
(Section~\ref{sec:engine_quality}); neither has a sealed-final result of its own. Latency was measured without the network and not under
concurrent load (Sections~\ref{sec:serving} and~\ref{sec:engine}); on TypeSafe's workflow
examples the served model is 3 to 19 times slower per case than the times TypeSafe published
for Jev (Section~\ref{sec:jev}).

\subsection{Operations: failures and reruns}
\label{sec:ops}

Appendix~\ref{app:compute} lists the compute behind each result. The following failures
occurred; none changed a reported result
except as stated.
\begin{itemize}
  \item \textbf{Evaluation v1.} The first readout ran on the flawed v1 evaluation
    (Section~\ref{sec:v1bug}); its results are kept separately and not used.
  \item \textbf{Backbone comparison.} A virtual environment that inherited system packages
    replaced Torch and crashed on import; the job was changed to use the image's environment
    and to refuse Torch replacement. An FP8 checkpoint tried to download kernels from the
    model hub; we declined to run third-party binaries on a host with storage permissions and
    used the BF16 weights. This exposed a bug that counted a loader error as 3,188 row
    failures; afterwards only over-length inputs count as row failures. Single 80\,GB GPUs
    were unavailable, so we used four-GPU L40S hosts. One large host idled for about 3.2
    hours after its job ended.
  \item \textbf{DeepSeek-V4.1-Flash.} A dependency check's regular expression mistook a new
    dependency for a Torch replacement and stopped the first attempt, after which the host
    shut down as designed. The second attempt hit a version conflict between two
    kernel libraries and a missing launcher on the path; both were fixed on the same host.
    36 rows returned NaN logits, of unknown cause, and count as failures.
  \item \textbf{Training.} We once judged the baseline process finished while it was still
    writing; results were unaffected. The convolution path ran without
    \texttt{causal\_conv1d}.
  \item \textbf{Sealed final and stress.} The first host finished and uploaded the final
    evaluation, then a post-processing error failed the job; the host powered itself off and
    the probe and benchmark results on its local disk were lost. The job script now
    uploads each result as soon as it exists and no longer fails on post-processing errors.
    The final evaluation was not rerun (the script refuses to rerun it when a final result
    exists); a second host reran only the probes, benchmarks and pointer retest. The first
    attempt to obtain that host failed for lack of capacity.
\end{itemize}
